\documentclass[12pt]{article}
\usepackage[utf8]{inputenc}
\usepackage[T1]{fontenc}
\usepackage{hyperref}
\title{Thai Sign Language Language Resources for K-12 Curriculum Development}
\author{Compiled by Claude Sonnet 5 (claude-sonnet-5) for Andr\'as Kornai}
\date{}
\begin{document}
\maketitle

\section{Introduction}
Thai Sign Language (ThSL; ISO 639-3 tsq; Glottocode thai1240) is the language of Thailand's deaf community. Its roots are in American Sign Language, introduced in deaf schools in the 1950s, blended with older local sign systems. A 1999 Education Ministry resolution recognised it as the national language of deaf people. Standardisation is still incomplete, and sources describe no single standard ThSL.

\section{Primary Educational Resources}
\begin{itemize}
\item Free online basic sign language course from Chiang Mai University on Thai MOOC, and free Saturday sessions at Setthasathian School, as listed on a Suan Sunandha Rajabhat University library page: \url{https://library.ssru.ac.th/news/view/s23126701}
\end{itemize}
No full K-12 curriculum or textbook series in ThSL was found.

\section{Linguistic and Reference Materials}
\begin{itemize}
\item Thai Sign Language Digital Dictionary project, Kasetsart University (preface, 2007): \url{https://pirun.ku.ac.th/~fhumalt/THSL/new_ThSL%20Digital%20Dict_23%20%20May%202007/Eng/preface.htm}
\item Another preface page of the Kasetsart ThSL dictionary: \url{https://pirun.ku.ac.th/~fhumalt/THSL/THSL/html/nav_en/THSL_pref_e.htm}
\item ``Wild sign diet'': an online sharing dictionary for the hearing impaired (KMUTT, 2011). It reports that experts favoured signs from Setsatian School for the Deaf: \url{https://kirim.kmutt.ac.th/converis/portal/detail/Publication/600529122?lang=en_GB}
\item Overview in Wikipedia, which cites the print \textit{Thai Sign Language Dictionary} (National Association of the Deaf in Thailand, 1986; revised 1990): \url{https://en.wikipedia.org/wiki/Thai_Sign_Language}
\end{itemize}

\section{Religious and Community Resources}
The National Association of the Deaf in Thailand published the print dictionary. No community website was found in the searches.

\section{Digital and Audio Resources}
The Kasetsart digital dictionary and the KMUTT online dictionary are the digital resources found. Their current availability is unconfirmed.

\section{Curriculum Development Context}
Most of what was found is a decade or more old. Standardisation is lacking, and the sources do not describe how the sign language is used in instruction. A curriculum project would need current information from Thai deaf schools such as Setthasathian, and from the National Association of the Deaf in Thailand.
\end{document}